\chapter{产品认证、环境可靠性与量产质量}
\label{chap:product-qualification-quality}

一块样板能够运行，不代表产品可以合法销售、稳定制造或在声明环境中持续工作。产品交付需要把法规适用性、EMC/安规、机械与热设计、环境验证、生产测量和现场失效反馈组织为闭环。

标准及法规具有地区、产品类别和版本边界。本章只给出通用工程方法；正式项目必须由合格实验室、认证机构和法规负责人确认适用条款。不得把“采用通过认证的模块”理解为整机自动合规。

\section{合规适用性矩阵}

立项阶段应建立矩阵，至少记录：销售地区、产品类别、供电方式、无线频段、使用环境、用户可接触接口、安装方式和预期寿命。每个适用法规或标准关联责任人、样机配置、测试实验室、计划时间和证据位置。

典型维度包括：
\begin{itemize}
  \item 电磁发射与抗扰度；
  \item 电气安全、绝缘、温升和可触及能量；
  \item 无线频谱、射频暴露和天线限制；
  \item 环境与机械适应性；
  \item 有害物质、材料和回收要求；
  \item 行业特定功能安全、网络安全或计量要求。
\end{itemize}

矩阵还应标识哪些变更会触发重测，例如时钟频率、DC/DC、PCB 层叠、外壳开孔、线缆、天线、无线固件、关键安全器件和电源适配器。

\section{EMC 发射}

传导与辐射发射来自差模电流、共模电流、时钟谐波、开关电源、线缆天线和回流路径。信息技术/多媒体设备常参考 CISPR 32 等产品族标准，具体限值、端口和测试距离取决于产品类别与地区\cite{cispr32,fcc-part15}。

设计阶段应把频率与物理结构关联：
\begin{itemize}
  \item 已知时钟、DDR、SerDes、PWM 和 DC/DC 的基频及谐波；
  \item 高速信号的回流路径、层切换和参考平面连续性；
  \item 外部线缆进入处的共模滤波、屏蔽连接和机壳地策略；
  \item 电源环路面积、开关节点铜皮和磁性器件漏磁；
  \item 软件扩频、驱动强度、边沿速率和工作模式的影响。
\end{itemize}

预扫应保存频谱、天线方向、线缆摆放、软件负载和样机配置。只保存“通过”照片无法支持整改比较。发现峰值后可以用近场探头、电流探头、分块停钟和替换线缆定位源与耦合路径。

\section{EMC 抗扰度}

常见抗扰度项目包括 ESD、EFT/burst、surge、射频场、传导射频、电压跌落和中断。IEC 61000-4 系列给出多类基础测试方法，例如 ESD、EFT 和 surge\cite{iec61000-4-2,iec61000-4-4,iec61000-4-5}；产品标准会进一步规定等级、端口和性能判据。

测试前必须定义性能判据：
\begin{description}
  \item[持续正常] 允许短暂数值扰动，但功能和通信不中断；
  \item[自动恢复] 允许可检测的重试或复位，规定恢复时间且不得丢失关键数据；
  \item[人工恢复] 仅在标准允许的项目中接受，并保证无危险状态；
  \item[不可接受] 锁死、错误执行器动作、配置损坏、安全边界失效或无法追踪的静默数据错误。
\end{description}

抗扰设计不只是增加 TVS。必须建立能量路径：连接器处泄放、机壳/保护地回流、信号限流、钳位电流、隔离边界、复位监控和软件恢复。ESD 后应用仍运行并不证明正确，应校验通信计数、传感器合理性、存储一致性和安全输出。

\section{电气安全与绝缘}

安规设计从工作电压、过电压类别、污染等级、材料组别、海拔和绝缘类型推导电气间隙、爬电距离与耐压要求。不能使用单一“安全距离”覆盖所有产品。

评审应检查：
\begin{itemize}
  \item 基本、附加、双重或加强绝缘的边界；
  \item 保险丝、限流、热熔断器和保护接地的失效模式；
  \item 用户可触及金属、接口和 SELV/危险电压边界；
  \item 变压器、光耦、继电器、Y 电容等安全关键器件认证；
  \item 单故障下温升、火灾、触电和机械危险；
  \item PCB 槽、阻焊和污染是否真的可计入绝缘设计。
\end{itemize}

安全关键 BOM 项目不应在采购阶段随意替代。替代料需要重新核对额定值、认证、封装距离、失效模式和温升。

\section{无线与天线变更}

采用预认证无线模块可以减少部分射频测试，但整机天线型号、增益、馈线、外壳、共存发射机和安装方式仍可能影响授权条件。欧盟无线设备指令等法规还覆盖安全、EMC 和频谱有效使用\cite{eu-red}。

量产测试不能只读模块序列号。应验证晶振偏差、发射功率、接收灵敏度代理指标、天线连接、校准数据和区域配置。固件必须防止加载错误地区的信道或功率参数。

\section{环境与机械验证}

IEC 60068 系列包含温度、湿热、振动、冲击等环境试验方法\cite{iec60068-series}。项目应根据真实运输、安装和使用环境选择应力谱，而不是机械地执行一组通用测试。

环境计划通常包括：
\begin{itemize}
  \item 高低温工作、存储与温度循环；
  \item 恒定/交变湿热与凝露风险；
  \item 随机振动、正弦扫频、机械冲击和跌落；
  \item 盐雾、腐蚀性气体、粉尘、进水和紫外老化；
  \item 连接器插拔、按键、继电器和风扇寿命；
  \item 包装运输与线束拉力。
\end{itemize}

试验应在应力期间运行代表性功能并记录连续数据。只在试验结束后开机，可能遗漏瞬时复位、接触抖动和传感器漂移。

\section{热设计与降额}

热设计要从环境温度、功耗分布、热界面、外壳、气流和安装方向建立路径。芯片结温可用简化热阻初估，但多热源、瞬态负载和封闭外壳需要仿真或实测修正。

验证至少覆盖：
\begin{itemize}
  \item 最大合法 CPU/GPU/NPU、网络、存储和显示并发负载；
  \item 风扇堵转、散热器接触退化和通风口遮挡；
  \item 温度传感器位置与真实热点之间的偏差；
  \item thermal throttling 对实时截止期和服务质量的影响；
  \item 热循环造成的焊点、连接器和密封应力。
\end{itemize}

器件降额应覆盖电压、电流、功耗、结温、纹波、机械应力和寿命。电解电容、风扇、继电器、Flash 和 LED 等寿命敏感器件需要依据实际应力建立维护或保修边界。

\section{可靠性模型与验证}

可靠性不能由一个 MTBF 数字概括。项目应区分随机硬件失效、系统性设计缺陷、磨损、环境诱发失效和制造缺陷。FMEA/FMEDA、故障树和降额分析用于发现薄弱环节，但结论需要试验与现场数据更新。

器件级 qualification 只能说明器件在规定条件下通过对应应力，不能代替整板互连、散热和软件工作负载验证。集成电路常使用 JEDEC 可靠性 qualification 方法，具体适用性取决于器件类型和供应商声明\cite{jedec-jesd47}。

可靠性指标应与任务剖面关联：每天启停次数、温度驻留、写入量、继电器动作、马达负载、振动时长和网络重连次数。没有任务剖面的寿命外推通常不可审计。

\section{HALT、HASS 与边界}

HALT 使用逐步升高的温度、振动和快速变化寻找设计裕量与薄弱环节，其目的通常是发现限制，不是直接证明寿命或法规合规。发现的失效应区分工作极限、破坏极限和夹具/试验设备问题。

HASS 或生产应力筛选用于暴露制造潜在缺陷，筛选强度必须建立在产品裕量和缺陷检出证据上。未经验证的过度筛选会消耗产品寿命或制造新缺陷。

任何加速试验都要说明加速模型、失效机理是否保持一致以及外推不确定度。高温不能加速所有机械或软件故障，组合应力也可能产生现场从未存在的新机理。

\section{EVT、DVT 与 PVT}

阶段名称因团队而异，但应具有明确退出证据：

\begin{description}
  \item[EVT] 验证关键架构、电源、时钟、接口、热路径和高风险器件；允许工程返修，但必须记录其对结论的影响。
  \item[DVT] 使用接近量产的硬件、固件和结构件关闭规格、环境、EMC、安规、可靠性和异常恢复。
  \item[PVT] 使用量产线、正式工装、生产软件、人员和供应链验证节拍、良率、追溯、返修与放行流程。
\end{description}

不能用“已经做了三轮板子”代替阶段门。若 DVT 样机仍有飞线、非量产散热器或工程电源，相关测试必须标记为条件性证据并在正式配置复验。

\section{测量不确定度与 GR\&R}

测试限值必须考虑 DUT 真实容差和测量系统不确定度。测量结果接近规格边界时，仪器精度、夹具压降、探针接触、环境、算法和重复性都会影响误判。测量不确定度可依据 GUM 方法组织来源与合成\cite{jcgm100}。

Gauge R\&R 用于评估重复性和再现性：同一操作员/设备重复测量的变化，以及不同操作员、工站或夹具之间的变化。自动工站也需要分析工装、参考源和接触差异，不能因为没有人工读数便假定再现性为零。

Guard band 策略应由风险决定。简单放宽生产限值会增加漏测，简单收紧会降低良率；需要结合测量能力、返修成本和现场风险选择。

\section{SPC、良率与缺陷 Pareto}

量产数据应保存连续测量值，而不仅是 Pass/Fail。控制图可以识别均值漂移、方差变化、工装老化和批次异常；控制界限描述过程行为，规格界限描述产品要求，二者不能混用。

建议持续跟踪：
\begin{itemize}
  \item 一次通过率、最终良率、返修率和报废率；
  \item 各测试项失败率和缺陷 Pareto；
  \item 关键模拟值、功耗、射频和校准参数分布；
  \item 工站间差异、夹具寿命和参考板趋势；
  \item 供应商、日期码、产线和班次相关性；
  \item 假失败与漏测的已知样本。
\end{itemize}

提高直通率的正确路径是消除缺陷或改善测量系统，而不是删除失败项、重复测试到通过或无审批修改限值。

\section{装联标准与外观判定}

焊接和装联外观应依据产品等级、图纸和适用标准形成统一判定。IPC-A-610 等文件提供电子组件可接受性框架，但客户规范和工艺文件可能提出额外要求\cite{ipc-a-610}。

外观缺陷与电气可靠性不是一一对应。BGA 空洞、底部焊点和内部裂纹可能需要 X-ray、切片、染色或电气应力；漂亮焊点也可能具有错误器件、错误程序或潜在污染。

\section{配置、替代料与变更控制}

量产基线应绑定 PCB/PCBA revision、BOM、AVL、固件、Bootloader、校准、测试程序、限值、工装和认证样机。任何单项独立升级都可能使证据失效。

替代料评审至少比较：
\begin{itemize}
  \item 引脚、封装、额定值和温度等级；
  \item 电气动态特性、启动默认值和时序；
  \item 软件 ID、寄存器、勘误和驱动兼容性；
  \item EMC、热、寿命、安规认证和供应状态；
  \item 需要重跑的设计验证、生产测试与法规项目。
\end{itemize}

PCN/EOL 应进入风险清单并保留 last-time-buy、重新设计和多来源策略。库存数量不能消除软件维护和认证重测周期。

\section{RMA、FRACAS 与 8D}

现场退回首先要保护证据：设备身份、版本、运行小时、环境、崩溃记录、存储状态和用户操作。直接重刷或返修可能永久破坏根因线索。

FRACAS 将 failure report、analysis 和 corrective action 关联起来。根因应区分：
\begin{itemize}
  \item 触发事件；
  \item 设计或制造根因；
  \item 为什么验证和生产测试没有发现；
  \item 为什么现场影响被放大；
  \item 临时遏制与永久纠正措施。
\end{itemize}

8D 等结构化方法适合跨团队问题，但模板不能代替物理证据。纠正措施完成应包括设计/工艺变化、验证结果、受影响批次处置，以及能防止同类缺陷再次逃逸的测试或规则。

现场故障率应按硬件版本、软件版本、批次、使用时长和任务剖面切片。单纯统计退货总数会把装机量差异和早期寿命效应混在一起。

\section{产品放行门}

产品放行至少应确认：
\begin{itemize}
  \item 法规适用性和所有证书/报告对应正式配置；
  \item 环境、热、机械和寿命验证覆盖声明任务剖面；
  \item DVT 缺陷、偏差和残余风险已有批准结论；
  \item PVT 良率、节拍、GR\&R、工装能力和追溯达到门槛；
  \item 安全关键与认证关键物料已经锁定；
  \item 现场诊断、RMA 保全证据和升级恢复流程可执行；
  \item 后续 PCN、法规变更、证书到期和可靠性趋势具有责任人。
\end{itemize}

\section{检查清单}

\begin{itemize}
  \item 是否依据销售地区、产品类别和使用环境建立合规矩阵？
  \item EMC 测试是否定义软件负载、线缆布局和功能性能判据？
  \item 环境与寿命试验是否源于真实任务剖面并在应力期间监测功能？
  \item EVT/DVT/PVT 是否具有配置一致的退出证据？
  \item 生产限值是否包含测量不确定度，GR\&R 和 SPC 是否持续运行？
  \item 变更和替代料是否评估认证、环境、软件和生产测试影响？
  \item RMA 是否保护证据并把根因转化为设计与测试改进？
\end{itemize}
